\documentclass[11pt]{article}
% ===========================================================================
%  Shared preamble for the AP Calculus AB exam-preparation worksheets.
%  Input right after \documentclass. Files beginning with "_" are not compiled
%  on their own. Navy and gold, to match the course's lesson pages.
% ===========================================================================
\usepackage[letterpaper,top=2.3cm,bottom=2.2cm,left=1.7cm,right=1.7cm,headheight=28pt,footskip=30pt]{geometry}
\usepackage{amsmath,amssymb}
\usepackage[T1]{fontenc}
\usepackage[utf8]{inputenc}
\usepackage{lmodern}
\usepackage{xcolor}
\usepackage[most]{tcolorbox}
\usepackage{enumitem}
\usepackage{fancyhdr}
\usepackage{array}
\usepackage{tikz}
\usepackage{pgfplots}
\pgfplotsset{compat=1.18}

\definecolor{navy}{HTML}{161E54}
\definecolor{ink}{HTML}{27357E}
\definecolor{gold}{HTML}{B98A00}
\definecolor{goldbg}{HTML}{FFFDF3}
\definecolor{inkbg}{HTML}{F6F7FD}
\definecolor{rule}{HTML}{C9CFF0}
\definecolor{okgreen}{HTML}{1B7A44}
\definecolor{okbg}{HTML}{F1F8F3}

\pagestyle{fancy}
\fancyhf{}
\fancyhead[L]{\footnotesize Name: \underline{\hspace{4cm}}}
\fancyhead[C]{\textbf{\color{navy}Integration Academy}\\[-2pt]{\scriptsize Calculus AB \textperiodcentered{} AP\textsuperscript{\textregistered} Exam Preparation}}
\fancyfoot[L]{\tiny\color{black!55}\parbox[t]{0.78\textwidth}{AP\textsuperscript{\textregistered} is a registered trademark of the College Board, which was not involved in the production of, and does not endorse, this material. \textcopyright{} Integration Academy.}}
\fancyfoot[R]{\scriptsize\thepage}
\renewcommand{\headrulewidth}{1.2pt}
\renewcommand{\headrule}{\hbox to\headwidth{\color{gold}\leaders\hrule height \headrulewidth\hfill}}
\renewcommand{\footrulewidth}{0pt}
\setlength{\parindent}{0pt}
\setlength{\parskip}{4pt}

% Title block: \wstitle{code}{title}{unit line}
\newcommand{\wstitle}[3]{\begin{center}{\small\bfseries\color{gold}WORKSHEET #1}\\[3pt]{\LARGE\bfseries\color{navy}#2}\\[3pt]{\small\color{black!60}#3}\end{center}\vspace{2pt}}
% Section heading
\newcommand{\sech}[1]{\par\vspace{10pt}{\large\bfseries\color{navy}#1}\par\nobreak\vspace{2pt}{\color{rule}\hrule height0.8pt}\vspace{7pt}}
% A level tag printed beside a group of questions
\newcommand{\level}[1]{\par\vspace{6pt}{\footnotesize\bfseries\color{gold}\MakeUppercase{#1}}\par\nobreak\vspace{2pt}}

\newtcolorbox{ideabox}{colframe=ink,colback=inkbg,coltitle=white,colbacktitle=ink,fonttitle=\bfseries\small,title={Key ideas},arc=2pt,boxrule=0.8pt,left=8pt,right=8pt}
\newtcolorbox{example}[1]{colframe=rule,colback=white,coltitle=navy,colbacktitle=inkbg,fonttitle=\bfseries\small,title={#1},arc=2pt,boxrule=0.8pt,borderline west={3pt}{0pt}{ink},left=10pt,right=8pt}
\newtcolorbox{frbox}[1]{colframe=gold,colback=goldbg,coltitle=white,colbacktitle=gold,fonttitle=\bfseries\small,title={#1},arc=2pt,boxrule=0.8pt,breakable,left=8pt,right=8pt}
\newtcolorbox{solbox}{colframe=okgreen,colback=okbg,coltitle=white,colbacktitle=okgreen,fonttitle=\bfseries\small,title={Worked solutions},arc=2pt,boxrule=0.8pt,breakable,left=8pt,right=8pt}
\newtcolorbox{scorebox}[1]{colframe=navy,colback=white,coltitle=white,colbacktitle=navy,fonttitle=\bfseries\small,title={#1},arc=2pt,boxrule=0.8pt,breakable,left=8pt,right=8pt}

\newcommand{\soln}{\par\smallskip\textbf{\color{ink}Solution.}\ }
\newcommand{\result}[1]{\par\smallskip\textbf{#1}}
% blank working space of a given height
\newcommand{\work}[1]{\par\vspace{#1}}
% a point value in the margin of a scoring guide
\newcommand{\pts}[1]{\hfill{\footnotesize\bfseries\color{gold}[#1]}}

% Coordinate grid for graph questions: \begin{apgraph}{xmin}{xmax}{ymin}{ymax} ... \end{apgraph}
\newenvironment{apgraph}[4]{\begin{center}\begin{tikzpicture}\begin{axis}[
  axis lines=middle, xmin=#1, xmax=#2, ymin=#3, ymax=#4,
  grid=both, grid style={black!12}, xtick distance=1, ytick distance=1,
  ticklabel style={font=\scriptsize}, xlabel={$x$}, ylabel={$y$},
  width=10cm, height=7.2cm, samples=80, clip=false,
  every axis plot/.append style={line width=1.1pt, color=ink}]}{\end{axis}\end{tikzpicture}\end{center}}
\newcommand{\closeddot}[2]{\addplot[only marks, mark=*, mark size=2.2pt, color=ink] coordinates {(#1,#2)};}
\newcommand{\opendot}[2]{\addplot[only marks, mark=*, mark size=2.2pt, color=ink, mark options={fill=white}] coordinates {(#1,#2)};}

% Multiple-choice options: four across, or stacked when the options are long.
\newcommand{\opts}[4]{\par\vspace{3pt}\begin{tabular}{@{}p{0.235\linewidth}p{0.235\linewidth}p{0.235\linewidth}p{0.235\linewidth}@{}}(A)\ #1 & (B)\ #2 & (C)\ #3 & (D)\ #4\end{tabular}\par\vspace{9pt}}
\newcommand{\optsv}[4]{\par\vspace{3pt}\begin{tabular}{@{}l@{\ \ }p{0.9\linewidth}@{}}(A) & #1\\(B) & #2\\(C) & #3\\(D) & #4\end{tabular}\par\vspace{9pt}}

\fancyhead[R]{\footnotesize\color{black!60}Worksheet 1.3}
\begin{document}

\wstitle{1.3}{Limits from Numerical Evidence}{Unit 1: Limits and Continuity}

\begin{ideabox}
\begin{itemize}[leftmargin=*,itemsep=2pt,topsep=2pt]
  \item Estimate $\displaystyle\lim_{x\to c}f(x)$ from a table by looking at inputs approaching $c$ from \textbf{both} sides.
  \item Read the \emph{trend} of the outputs, not the single entry nearest to $c$.
  \item A table is evidence, not proof: write \textquotedblleft{}the table suggests\ldots\textquotedblright{}.
  \item For trigonometric functions the calculator must be in \textbf{radian} mode.
\end{itemize}
\end{ideabox}

\sech{Worked examples}
\begin{example}{Example 1 \textperiodcentered{} Estimating a limit}
Selected values of $f$ are shown. Estimate $\displaystyle\lim_{x\to 4}f(x)$.\begin{center}\renewcommand{\arraystretch}{1.25}\begin{tabular}{|c|c|c|c|c|}\hline $x$ & $3.9$ & $3.99$ & $4.01$ & $4.1$ \\ \hline $f(x)$ & $1.9$ & $1.99$ & $2.01$ & $2.1$ \\ \hline\end{tabular}\end{center}\soln From the left the values rise toward 2; from the right they fall toward 2.\result{The table suggests the limit is $2$.}
\end{example}

\begin{example}{Example 2 \textperiodcentered{} Sides that disagree}
Selected values of $g$ are shown. What does the table suggest about $\displaystyle\lim_{x\to 1}g(x)$?\begin{center}\renewcommand{\arraystretch}{1.25}\begin{tabular}{|c|c|c|c|c|}\hline $x$ & $0.9$ & $0.99$ & $1.01$ & $1.1$ \\ \hline $g(x)$ & $-1.2$ & $-1.02$ & $3.03$ & $3.3$ \\ \hline\end{tabular}\end{center}\soln Left-hand values approach $-1$; right-hand values approach $3$.\result{The table suggests that the two-sided limit does not exist.}
\end{example}

\begin{example}{Example 3 \textperiodcentered{} Building a table (calculator)}
Estimate $\displaystyle\lim_{x\to 0}\dfrac{e^{3x}-1}{x}$ to three decimal places.\soln At $x=0.001$ the value is $3.0045$; at $x=-0.001$ it is $2.9955$. At $x=\pm0.0001$ the values are $3.00045$ and $2.99955$.\result{The limit is approximately $3.000$.}
\end{example}

\newpage
\sech{Practice}

\level{Core}
\begin{enumerate}[leftmargin=*,itemsep=2pt]
  \item A table gives $f(5.9)=11.8$, $f(5.99)=11.98$, $f(6.01)=12.02$, $f(6.1)=12.2$. Estimate $\displaystyle\lim_{x\to 6}f(x)$.\work{2.3cm}
  \item A table gives $h(1.9)=6.7$, $h(1.99)=6.97$, $h(2)=-4$, $h(2.01)=7.03$, $h(2.1)=7.3$. Estimate $\displaystyle\lim_{x\to 2}h(x)$.\work{2.3cm}
  \item A table gives values of $k$ only at $x=2.9,\ 2.99,\ 2.999$, which approach 8. What can and cannot be concluded?\work{2.3cm}
  \item Which inputs would you choose to investigate $\displaystyle\lim_{x\to -3}f(x)$? Give six.\work{2.3cm}
\end{enumerate}

\level{Exam level}
\begin{enumerate}[leftmargin=*,itemsep=2pt,start=5]
  \item (Calculator) Estimate $\displaystyle\lim_{x\to 0}\dfrac{\sin(2x)}{\sin(5x)}$ to three decimal places.\work{2.8cm}
  \item (Calculator) Estimate $\displaystyle\lim_{x\to 1}\dfrac{x^5-1}{x-1}$.\work{2.8cm}
  \item A student finds $\dfrac{\sin x}{x}\approx0.01745$ at $x=0.001$. What has gone wrong, and what is the correct value?\work{2.8cm}
  \item A function $f$ is defined for all real numbers. A table shows $f(2.9)=5.8$, $f(2.99)=5.98$, $f(3.01)=6.02$ and $f(3.1)=6.2$. Must $\displaystyle\lim_{x\to 3}f(x)=6$? Explain.\work{2.8cm}
\end{enumerate}

\begin{frbox}{Question 9 \textperiodcentered{} Free-response practice}
(Calculator) Let $h(x)=\dfrac{\sqrt{x+1}-1}{x}$ for $x\neq0$.
\begin{enumerate}[label=(\alph*),leftmargin=*,itemsep=2pt]
  \item Find $h(-0.1)$, $h(-0.01)$, $h(0.01)$ and $h(0.1)$ to four decimal places.\work{2.6cm}
  \item Use your values to estimate $\displaystyle\lim_{x\to 0}h(x)$.\work{2.6cm}
  \item Explain why part (b) is not a proof, and confirm the limit algebraically.\work{2.6cm}
\end{enumerate}
\end{frbox}

\newpage
\begin{solbox}
\begin{enumerate}[leftmargin=*,itemsep=4pt]
  \item $12$.
  \item $7$. The entry $h(2)=-4$ is the function value, not the limit.
  \item The table suggests $\displaystyle\lim_{x\to 3^-}k(x)=8$. It gives no information about the right-hand limit, so nothing can be said about the two-sided limit.
  \item For example $-3.1,\ -3.01,\ -3.001,\ -2.999,\ -2.99,\ -2.9$: three on each side, closing in.
  \item At $x=\pm0.01$ the value is $0.40014$; at $\pm0.001$ it is $0.4000014$. The limit is about $0.400$.
  \item At $x=1.001$: $5.0100$. At $x=0.999$: $4.9900$. The limit is about $5$.
  \item The calculator is in degree mode ($0.01745\approx\tfrac{\pi}{180}$). In radian mode the value is $0.9999998$, close to 1.
  \item No. Four values suggest the limit is 6 but cannot force it; $f$ is unrestricted at inputs closer to 3 than those listed.
  \item \begin{enumerate}[label=(\alph*),leftmargin=*,itemsep=3pt]
    \item $h(-0.1)=0.5132$, $h(-0.01)=0.5013$, $h(0.01)=0.4988$, $h(0.1)=0.4881$.
    \item From both sides the values approach $0.5$.
    \item A table samples only finitely many inputs. Algebraically, $\dfrac{\sqrt{x+1}-1}{x}\cdot\dfrac{\sqrt{x+1}+1}{\sqrt{x+1}+1}=\dfrac{1}{\sqrt{x+1}+1}\to\dfrac12$.
  \end{enumerate}
\end{enumerate}
\end{solbox}

\end{document}
